\chapter{Summary of results}\label{ch:summary}
This chapter lists the results of Chapters~\ref{ch:I}--\ref{ch:W} with their status labels (\S\ref{gl:labels}) and pointers, one section per chapter.

\begingroup
% macros of Part~I, Draft~11
\providecommand{\ov}{}\renewcommand{\ov}[1]{\overline{#1}}
\providecommand{\eps}{}\renewcommand{\eps}{\epsilon}
\providecommand{\A}{}\renewcommand{\A}{\mathbb{A}}
\providecommand{\C}{}\renewcommand{\C}{\mathbb{C}}
\providecommand{\Q}{}\renewcommand{\Q}{\mathbb{Q}}
\providecommand{\Z}{}\renewcommand{\Z}{\mathbb{Z}}
\providecommand{\cO}{}\renewcommand{\cO}{\mathcal{O}}
\providecommand{\gen}{}\renewcommand{\gen}{\mathrm{gen}}
\providecommand{\reg}{}\renewcommand{\reg}{\mathrm{reg}}
\providecommand{\tr}{}\renewcommand{\tr}{\operatorname{tr}}
\providecommand{\Gal}{}\renewcommand{\Gal}{\operatorname{Gal}}
\providecommand{\Aut}{}\renewcommand{\Aut}{\operatorname{Aut}}
\providecommand{\Nrd}{}\renewcommand{\Nrd}{\operatorname{Nrd}}
\providecommand{\Ad}{}\renewcommand{\Ad}{\operatorname{Ad}}
\providecommand{\Hom}{}\renewcommand{\Hom}{\operatorname{Hom}}
\providecommand{\Iso}{}\renewcommand{\Iso}{\operatorname{Iso}}
\providecommand{\vol}{}\renewcommand{\vol}{\operatorname{vol}}
\providecommand{\supp}{}\renewcommand{\supp}{\operatorname{supp}}
\providecommand{\diag}{}\renewcommand{\diag}{\operatorname{diag}}
\providecommand{\GL}{}\renewcommand{\GL}{\mathrm{GL}}
\providecommand{\SL}{}\renewcommand{\SL}{\mathrm{SL}}
\providecommand{\Gb}{}\renewcommand{\Gb}{\ov{G}}
\providecommand{\Zb}{}\renewcommand{\Zb}{\ov{Z}}
\providecommand{\GDb}{}\renewcommand{\GDb}{\ov{G}_D}
\providecommand{\pib}{}\renewcommand{\pib}{\ov{\pi}}
\providecommand{\ks}{}\renewcommand{\ks}{\kappa^{*}}
\providecommand{\Eb}{}\renewcommand{\Eb}{\ov{E}}
\providecommand{\Tb}{}\renewcommand{\Tb}{\ov{T}}
\providecommand{\cG}{}\renewcommand{\cG}{\mathcal{G}}
\providecommand{\cGb}{}\renewcommand{\cGb}{\ov{\mathcal{G}}}
\providecommand{\iD}{}\renewcommand{\iD}{\iota_D}
\providecommand{\iG}{}\renewcommand{\iG}{\iota_{\GL}}
\providecommand{\jj}{}\renewcommand{\jj}{\jmath}
\providecommand{\Pv}{}\renewcommand{\Pv}{\textbf{(P)}}
\providecommand{\Ev}{}\renewcommand{\Ev}{\textbf{(E)}}
\providecommand{\Ov}{}\renewcommand{\Ov}{\textbf{(O)}}
\providecommand{\Cv}{}\renewcommand{\Cv}{\textbf{(C)}}
% =====================================================================
\section{\texorpdfstring{Local matching and a global transfer under elliptic-support conditions (Chapter~\ref{ch:I})}{Local matching and a global transfer under elliptic-support conditions (Chapter 4)}}\label{I:sec:summary}
The results of Chapter~\ref{ch:I}, with status labels (\S\ref{gl:labels}) and pointers; notation as in Assumption~\ref{I:ass:standing} and \S\ref{I:sec:glossary}, with $v\in S$ and $F=K_v$.

\begin{enumerate}[leftmargin=*,label=\arabic*.]
\item \textbf{One ambient cover} (Lemma~\ref{I:lem:center}, Proposition~\ref{I:prop:restriction}; \Pv\ using \Ev). The cover $\GDb(K_v)$ of \cite{JC26} and the Kazhdan--Patterson cover $\Gb_v$ of $\GL(r,K_v)$ are subgroups of one cover $\cGb_v$ of $\GL(r,L_v)$, with the same central subgroup $\Zb_v$. As an extension, $\Gb_v$ is the parameter-$c_v$ cover after relabeling $\mu_n$, and the parameter is forced (Remark~\ref{I:rem:forced}). At $v\notin S$ they are identified canonically by conjugation (Proposition~\ref{I:prop:outsideS}; \Pv).
\item \textbf{Theorem~A: covered elliptic tori} (Proposition~\ref{I:prop:comm}, Proposition~\ref{I:prop:good}, Theorem~\ref{I:thm:ident}, Corollaries~\ref{I:cor:classes} and~\ref{I:cor:reps}; \Pv). For every maximal subfield $E\subset D_v$ the two covered tori have the commutator pairing $((Nx)^{1+2c_v}x^{-1},y)_E$ and the good elements $p^{-1}(F^\times E^{\times n})$, of index $|n|_F^{-(r-1)}$. Every isomorphism between them over $\iG\circ j\circ\iD^{-1}$ that is the identity on $\Zb_v$ has the same restriction $\jj_{E,j}$ to the good elements. This gives a canonical bijection $\jj$ of good regular elliptic classes and of genuine representations of covered tori.
\item \textbf{Theorem~B: matching on the regular set} (Theorem~\ref{I:thm:easy}, \S\ref{I:app:oi}; \Pv). Matching (Definition~\ref{I:def:matching}) carries the sign $(-1)^{r-1}$. Every genuine $\tilde f_D$ vanishing near $\Zb_v$ has a matching $\tilde f$ supported on the good regular elliptic set, and conversely every regular-supported $\tilde f$ with vanishing non-elliptic orbital integrals has a matching $\tilde f_D$. Near the center the method fails (Proposition~\ref{I:prop:germ}; \Pv).
\item \textbf{Rational elements and trace formulas} (Theorem~\ref{I:thm:rational}, Corollary~\ref{I:cor:rationalbij}, Propositions~\ref{I:prop:geomD} and~\ref{I:prop:FKvariant}, Theorem~\ref{I:thm:comparison}; \Pv\ using the \Ev\ inputs of Remark~\ref{I:rem:Einputs}). The elliptic terms of the two trace formulas agree term by term; on the $\GL$ side we use a variant of \cite[\S18]{FK86} with the supercusp form at a place $w_0\notin S$.
\item \textbf{Packet identity} (Proposition~\ref{I:prop:packet}; \Pv, unconditional). The $D$-side and cuspidal $\GL$-side near-equivalence packets with the same components outside $S\cup\{w_0\}$ have equal traces at $S$ against matching regular-supported test functions; both sums are locally finite.
\item \textbf{Theorem~C: global transfer} (Theorem~\ref{I:thm:transfer}; \Pv\ under $(\mathrm{Hyp}_D)$). A genuine automorphic $\pib_D$ with a supercuspidal component at some finite $w_0\notin S$ transfers to a genuine cuspidal $\pib$ with the same components outside $S$, if its packet satisfies $(\mathrm{Hyp}_D)$; this holds when the packet is finite (Proposition~\ref{I:prop:HDfinite}; \Pv). The converse holds under (Hyp) (Theorem~\ref{I:thm:converse}), and Lemma~\ref{I:lem:cancel} (\Pv) shows exactly how $(\mathrm{Hyp}_D)$ could fail.
\item \textbf{What is left} (\S\ref{I:sec:deferred}; \Cv). (D1) matching of arbitrary functions, (D2) local character relations, (D3) finiteness of packets, (D4) composition with \cite{JC26}. Chapter~\ref{ch:II} proves (D1), and Chapter~\ref{ch:III} proves (D2) at tame places in the range of the corrected Flicker--Kazhdan correspondence (Remark~\ref{I:rem:D2tame}); the rest is open.
\end{enumerate}
\emph{Running example} (Example~\ref{I:ex:r2n3}): the covered elliptic tori are non-abelian at $\mathfrak p_3$ (index $9$, dimension $3$) and abelian at $\mathfrak p_7$; the sign in Theorem~B is $-1$.
\endgroup
\begingroup
% macros of Part~II, Draft~10
\providecommand{\ov}{}\renewcommand{\ov}[1]{\overline{#1}}
\providecommand{\A}{}\renewcommand{\A}{\mathbb{A}}
\providecommand{\C}{}\renewcommand{\C}{\mathbb{C}}
\providecommand{\Z}{}\renewcommand{\Z}{\mathbb{Z}}
\providecommand{\cO}{}\renewcommand{\cO}{\mathcal{O}}
\providecommand{\fp}{}\renewcommand{\fp}{\mathfrak{p}}
\providecommand{\gen}{}\renewcommand{\gen}{\mathrm{gen}}
\providecommand{\reg}{}\renewcommand{\reg}{\mathrm{reg}}
\providecommand{\tr}{}\renewcommand{\tr}{\operatorname{tr}}
\providecommand{\Nrd}{}\renewcommand{\Nrd}{\operatorname{Nrd}}
\providecommand{\Ad}{}\renewcommand{\Ad}{\operatorname{Ad}}
\providecommand{\Iso}{}\renewcommand{\Iso}{\operatorname{Iso}}
\providecommand{\vol}{}\renewcommand{\vol}{\operatorname{vol}}
\providecommand{\supp}{}\renewcommand{\supp}{\operatorname{supp}}
\providecommand{\Cent}{}\renewcommand{\Cent}{\operatorname{Cent}}
\providecommand{\GL}{}\renewcommand{\GL}{\mathrm{GL}}
\providecommand{\St}{}\renewcommand{\St}{\mathrm{St}}
\providecommand{\Gb}{}\renewcommand{\Gb}{\ov{G}}
\providecommand{\Zb}{}\renewcommand{\Zb}{\ov{Z}}
\providecommand{\GDb}{}\renewcommand{\GDb}{\ov{G}_D}
\providecommand{\pib}{}\renewcommand{\pib}{\ov{\pi}}
\providecommand{\Eb}{}\renewcommand{\Eb}{\ov{E}}
\providecommand{\Tb}{}\renewcommand{\Tb}{\ov{T}}
\providecommand{\cG}{}\renewcommand{\cG}{\mathcal{G}}
\providecommand{\cGb}{}\renewcommand{\cGb}{\ov{\mathcal{G}}}
\providecommand{\iD}{}\renewcommand{\iD}{\iota_D}
\providecommand{\iG}{}\renewcommand{\iG}{\iota_{\GL}}
\providecommand{\jj}{}\renewcommand{\jj}{\jmath}
\providecommand{\sg}{}\renewcommand{\sg}{\sigma}
\providecommand{\Pv}{}\renewcommand{\Pv}{\textbf{(P)}}
\providecommand{\Ev}{}\renewcommand{\Ev}{\textbf{(E)}}
\providecommand{\Ov}{}\renewcommand{\Ov}{\textbf{(O)}}
\providecommand{\Cv}{}\renewcommand{\Cv}{\textbf{(C)}}
% =====================================================================
\section{\texorpdfstring{Matching of arbitrary functions and an improved global transfer (Chapter~\ref{ch:II})}{Matching of arbitrary functions and an improved global transfer (Chapter 5)}}\label{II:sec:summary}
The results of Chapter~\ref{ch:II}, with status labels (\S\ref{gl:labels}) and pointers; notation as in \S\ref{II:sec:recall} and \S\ref{II:sec:glossary}. ``Relative to Chapter~\ref{ch:I}'' means that the cited results of Chapter~\ref{ch:I}, restated in \S\ref{II:subsec:fromPartI}, are used.

\begin{enumerate}[leftmargin=*,label=\arabic*.]
\item \textbf{Compatible measures} (Definition~\ref{II:def:measures}, Remark~\ref{II:rem:measures}; \Pv). At $v\in S$, $\vol(D_v^\times/F^\times)=d(\St)^{-1}$; every result of Chapter~\ref{ch:I} holds unchanged with these measures.
\item \textbf{Theorem~D: the section $\sg$ and transport} (Lemma~\ref{II:lem:roots}, Propositions~\ref{II:prop:section} and~\ref{II:prop:transport}; \Pv). At every finite place, including those above $n$, and on both covers: a canonical $n$-th root on a conjugation-invariant neighborhood $V_M$ of the unipotent set gives a section $\sg(g)=\widetilde{g^{1/n}}^{\,n}$ that is equivariant, good at regular points, compatible with the splittings and with $\jj$. The map $\mathcal T$ identifies genuine functions supported on $p^{-1}(ZV_M)$ with functions on $G/Z$ supported on $ZV_M$, preserving orbital integrals.
\item \textbf{Germs at central good elements} (Corollary~\ref{II:cor:germexp}; \Pv\ using \Ev). They are the Shalika germs of $\GL(r,F)$; on the $D$ side only the trivial germ occurs.
\item \textbf{Orbital integrals of arbitrary genuine functions} (Lemmas~\ref{II:lem:singular} and~\ref{II:lem:Aext}; \Pv). Away from $\Zb$ they are admissible in the sense of \S\ref{I:subsec:transported} ($r$ prime).
\item \textbf{Linear inputs and the Steinberg-type pair} (\S\ref{II:subsec:linear}, Proposition~\ref{II:prop:eD}). (L1) is \Ev; (L2) is \Pv\ using \Ev\ (Lemma~\ref{II:lem:L2}); $f_{\St}(1)=d(\St)$ (Lemma~\ref{II:lem:fSt1}; \Pv). The genuine pair $(e_D,\tilde f^{\St})$ matches (\Pv), but $\tilde f^{\St}$ is not a pseudo-coefficient of a single genuine representation (Remark~\ref{II:rem:notpc}).
\item \textbf{Theorem~E: matching of arbitrary functions} (Theorem~\ref{II:thm:matching}; \Pv\ using (L1), (L2)). Every genuine $\tilde f_D$ has a matching $\tilde f=\tilde f_D(1)\tilde f^{\St}+\tilde f_{\reg}$, and every genuine $\tilde f$ with vanishing non-elliptic orbital integrals has a matching $\tilde f_D$.
\item \textbf{The required germ} (Theorem~\ref{II:thm:germ}; \Pv; Lemma~\ref{II:lem:cuspnear1} and Corollary~\ref{II:cor:trivgerm}, \Pv\ using \Ev). Near $\Zb_v$ a function matches $\tilde f_D$ exactly when its non-trivial unipotent orbital integrals vanish and $\tilde f(1)=\tilde f_D(1)$.
\item \textbf{Theorem~C$'$} (Theorem~\ref{II:thm:comparison2}, Proposition~\ref{II:prop:packet2}; \Pv\ relative to Chapter~\ref{ch:I}; Theorem~\ref{II:thm:Cprime}, \Pv\ under $(\mathrm{Hyp}_D')$). The packet identity holds with arbitrary matching functions at all places of $S$ but one, and $\pib_D$ transfers under $(\mathrm{Hyp}_D')$, which is implied by $(\mathrm{Hyp}_D)$. The converse holds under $(\mathrm{Hyp}')$ (Theorem~\ref{II:thm:Cprimeconv}), in particular for a single-member packet with square-integrable components at $S$ (Corollary~\ref{II:cor:Hprimesq}; \Pv\ using \Ev).
\item \textbf{When $(\mathrm{Hyp}_D')$ fails} (Proposition~\ref{II:prop:cancel2}, Corollary~\ref{II:cor:finiteStypes}; \Pv). Exactly when the $S$-multiplicities are proportional to the Plancherel weights; it holds if the levels at $S$ are bounded.
\item \textbf{What is left} (\S\ref{II:sec:partIII}; \Cv). Local character relations (since proved by Chapter~\ref{ch:III} at tame places in the range of the corrected Flicker--Kazhdan correspondence), an unconditional global correspondence, removal of $(\mathrm{Hyp}_D')$ and $(\mathrm{Hyp}')$, and completing \cite[Remark~10.3]{JC26}.
\end{enumerate}
\emph{Running example:} Theorem~E applies at both $\mathfrak p_3$ and $\mathfrak p_7$; at the wild place $\mathfrak p_3$ matching near the center uses Theorem~D.
\endgroup
\begingroup
% macros of Part~III, Draft~10
\providecommand{\A}{}\renewcommand{\A}{\mathbb A}
\providecommand{\C}{}\renewcommand{\C}{\mathbb C}
\providecommand{\Z}{}\renewcommand{\Z}{\mathbb Z}
\providecommand{\GL}{}\renewcommand{\GL}{\operatorname{GL}}
\providecommand{\Nrd}{}\renewcommand{\Nrd}{\operatorname{Nrd}}
\providecommand{\Irr}{}\renewcommand{\Irr}{\operatorname{Irr}}
\providecommand{\vol}{}\renewcommand{\vol}{\operatorname{vol}}
\providecommand{\tH}{}\renewcommand{\tH}{\widetilde H}
\providecommand{\tZ}{}\renewcommand{\tZ}{\widetilde Z}
\providecommand{\ttT}{}\renewcommand{\ttT}{\widetilde T}
\providecommand{\reg}{}\renewcommand{\reg}{\mathrm{reg}}
\providecommand{\gen}{}\renewcommand{\gen}{\mathrm{gen}}
\providecommand{\JL}{}\renewcommand{\JL}{\mathrm{JL}}
\providecommand{\Sha}{}\renewcommand{\Sha}{\operatorname{Sha}}
\providecommand{\Gb}{}\renewcommand{\Gb}{\overline{G}}
\providecommand{\GDb}{}\renewcommand{\GDb}{\overline{G}_D}
\providecommand{\Zb}{}\renewcommand{\Zb}{\overline{Z}}
\providecommand{\Ptag}{}\renewcommand{\Ptag}{\textnormal{\textbf{(P)}}}
\providecommand{\Etag}{}\renewcommand{\Etag}{\textnormal{\textbf{(E)}}}
\providecommand{\Otag}{}\renewcommand{\Otag}{\textnormal{\textbf{(O)}}}
\providecommand{\Ctag}{}\renewcommand{\Ctag}{\textnormal{\textbf{(C)}}}
\section{\texorpdfstring{The division-algebra side and its local character transform (Chapter~\ref{ch:III})}{The division-algebra side and its local character transform (Chapter 6)}}\label{III:sec:summary}
The results of Chapter~\ref{ch:III}, with status labels (\S\ref{gl:labels}) and pointers. Throughout, $F$ is a $p$-adic field containing $\mu_n$, $D/F$ is a central division algebra of prime degree $r$ with $(n,r)=1$, $H=D^\times$, and $\tH$ is the cover of Chapter~\ref{ch:I}.

\begin{enumerate}[leftmargin=*,label=\arabic*.]
\item \textbf{Good elements} (Proposition~\ref{III:prop:good}; \Etag\ from Chapter~\ref{ch:I}). For each maximal subfield $E$, the good elements of $p^{-1}(E^\times)$ are the $z\,t^*$, $z\in\tZ$, $t\in E^\times$.
\item \textbf{Power maps} (Lemmas~\ref{III:lem:regular}, \ref{III:lem:index}, \ref{III:lem:ratio}; \Ptag). The $n$-th power map creates no new singular elements; its degree on elliptic tori and the discriminant ratio $\Delta(t)/\Delta(t^n)$ are computed, including at wild places.
\item \textbf{The transform $L_D$} (Definition~\ref{III:def:L}, Lemma~\ref{III:lem:well-defined}, Theorem~\ref{III:thm:unitary}; \Ptag). $L_D$ is a unitary isomorphism between $\theta$-equivariant invariant $L^2$ functions on $H$ and genuine $\tilde\theta$-equivariant ones on $\tH$, with coefficient $b_D=|n|_F^{-(r-1)/2}$ forced by unitarity.
\item \textbf{The character matrix} (Corollaries~\ref{III:cor:matrix}, \ref{III:cor:integrality}, Lemma~\ref{III:lem:integral-constants}; \Ptag). The matrix of $L_D$ in irreducible characters is unitary; if it is integral, then $\dim\tau=c_D\dim\rho$ with $c_D=|n|_F^{-(r^2-1)/2}$; $b_D$ and $c_D$ are positive integers.
\item \textbf{Adams operations} (Proposition~\ref{III:prop:adams}; \Ptag). The Adams pullback and its weight, the bridge used at tame places.
\item \textbf{Tame places} (Lemma~\ref{III:lem:prime-to-n}, Proposition~\ref{III:prop:tame-split}, Theorem~\ref{III:thm:tame-transfer}; \Ptag). If $p\nmid n$, the cover splits uniquely over the norm-one group, and $L_D$ maps irreducible characters bijectively to genuine irreducible characters, preserving dimensions.
\item \textbf{The character relation (D2) at tame places} (Corollary~\ref{III:cor:D2tame}; \Ptag\ in the stated range of the corrected Flicker--Kazhdan correspondence, using \Etag). Composing $L_D$, the linear Jacquet--Langlands correspondence and Zou's corrected correspondence gives (D2) of Chapter~\ref{ch:I} with the sign $(-1)^{r-1}$.
\item \textbf{Wild places} (Proposition~\ref{III:prop:wild-obstruction}; \Ptag; Remark~\ref{III:rem:wild-test}). If $p\mid n$, the cover does not split over the norm-one group and genuine irreducibles have dimension divisible by $|n|_F^{-(r-1)/2}$. Integrality of $L_D$ there is \Ctag.
\item \textbf{Orbital-integral transfer} (Theorem~\ref{III:thm:test-transfer}, Proposition~\ref{III:prop:trace-pair}, Corollary~\ref{III:cor:identity-factor}; \Ptag). A transfer of invariant test functions at every place of $S$, adjoint to $L_D$ under the trace pairing.
\item \textbf{Global powers} (Propositions~\ref{III:prop:global-cokernel}, \ref{III:prop:global-power-hasse}, Corollary~\ref{III:cor:sha-power}; \Ptag). A local-to-global principle for $n$-th powers modulo scalars outside any finite set of places, and the vanishing of the associated $\Sha^1$. Section~\ref{III:sec:kottwitz} reviews Kottwitz's framework (\Etag).
\item \textbf{The global comparison} (\S\ref{III:sec:global}; \Ctag). The spectral identity \eqref{III:eq:global-conj} is a target, with four proof obligations.
\end{enumerate}
\emph{Running example:} at $\mathfrak p_7$, $b_D=c_D=1$ and $L_D$ is a dimension-preserving bijection; at $\mathfrak p_3$, $b_D=3$, $c_D=27$, and integrality is open here (Chapter~\ref{ch:IV} and the companion chapter settle the trivial character and all rows trivial on the third congruence subgroup).
\endgroup
\begingroup
% macros of Part~IV, concluding Draft~6
\providecommand{\F}{}\renewcommand{\F}{F}
\providecommand{\K}{}\renewcommand{\K}{K}
\providecommand{\A}{}\renewcommand{\A}{\mathbb A}
\providecommand{\C}{}\renewcommand{\C}{\mathbb C}
\providecommand{\Z}{}\renewcommand{\Z}{\mathbb Z}
\providecommand{\GL}{}\renewcommand{\GL}{\operatorname{GL}}
\providecommand{\Irr}{}\renewcommand{\Irr}{\operatorname{Irr}}
\providecommand{\Nrd}{}\renewcommand{\Nrd}{\operatorname{Nrd}}
\providecommand{\Nm}{}\renewcommand{\Nm}{\operatorname{N}}
\providecommand{\vol}{}\renewcommand{\vol}{\operatorname{vol}}
\providecommand{\Gal}{}\renewcommand{\Gal}{\operatorname{Gal}}
\providecommand{\Aut}{}\renewcommand{\Aut}{\operatorname{Aut}}
\providecommand{\reg}{}\renewcommand{\reg}{\mathrm{reg}}
\providecommand{\tH}{}\renewcommand{\tH}{\widetilde H}
\providecommand{\tZ}{}\renewcommand{\tZ}{\widetilde Z}
\providecommand{\tT}{}\renewcommand{\tT}{\widetilde T}
\providecommand{\Gb}{}\renewcommand{\Gb}{\overline G}
\providecommand{\GDb}{}\renewcommand{\GDb}{\overline{G_D}}
\providecommand{\Ptag}{}\renewcommand{\Ptag}{\textnormal{\textbf{(P)}}}
\providecommand{\Etag}{}\renewcommand{\Etag}{\textnormal{\textbf{(E)}}}
\providecommand{\Ctag}{}\renewcommand{\Ctag}{\textnormal{\textbf{(C)}}}
\providecommand{\Otag}{}\renewcommand{\Otag}{\textnormal{\textbf{(O)}}}
\section{\texorpdfstring{Local power maps, toral positivity, and conditional global comparison (Chapter~\ref{ch:IV})}{Local power maps, toral positivity, and conditional global comparison (Chapter 7)}}\label{IV:sec:summary}
We conclude the fourth part of the metaplectic Jacquet--Langlands
series with a precise separation of proved local and arithmetic
results from two still-conditional goals.  We audit the ambient
cocycle and scalar-central parameter, recall the all-torus
good-element criterion, compute kernels and Haar moduli of
$n$th-power maps on finite etale tori modulo scalars, and
classify globally good regular rational classes.  Building on
Chapter~\ref{ch:III}'s unitary division-side character transform, we prove
that wild irreducible-character integrality is equivalent to
positive definiteness of every transferred irreducible character.
A new toral permutation-character theorem establishes positive
integral Fourier coefficients of the discriminant ratio for
\emph{every quadratic field extension} over a nonarchimedean
field of odd residue characteristic, for all odd $n$ including
$p\mid n$.  The complete numerical $r=2$, $n=3$ example and the
nonabelian gluing problem are developed separately in a
standalone companion chapter (Chapter~\ref{ch:W}), here begun as Draft~0.  We compare
our normalizations with Zou's correction of Flicker--Kazhdan,
state explicit conditions for split-place matching, and prove
a conditional geometric and trace-distribution comparison with
normalization factor $\Lambda$.  No full wild-place irreducible
correspondence, construction of exceptional local matching
functions, or unconditional ordinary-to-covered global
comparison is claimed.

The status of each statement is tabulated in \S\ref{IV:sec:intro}.
\endgroup
\begingroup
% macros of the companion chapter on the rank-two wild place, Draft~8
\providecommand{\F}{}\renewcommand{\F}{F}
\providecommand{\Z}{}\renewcommand{\Z}{\mathbb Z}
\providecommand{\C}{}\renewcommand{\C}{\mathbb C}
\providecommand{\Q}{}\renewcommand{\Q}{\mathbb Q}
\providecommand{\Fq}{}\renewcommand{\Fq}{\mathbb F}
\providecommand{\Nm}{}\renewcommand{\Nm}{\operatorname N}
\providecommand{\Nrd}{}\renewcommand{\Nrd}{\operatorname{Nrd}}
\providecommand{\Trd}{}\renewcommand{\Trd}{\operatorname{Trd}}
\providecommand{\Tr}{}\renewcommand{\Tr}{\operatorname{Tr}}
\providecommand{\Irr}{}\renewcommand{\Irr}{\operatorname{Irr}}
\providecommand{\Ind}{}\renewcommand{\Ind}{\operatorname{Ind}}
\providecommand{\Ad}{}\renewcommand{\Ad}{\operatorname{Ad}}
\providecommand{\tH}{}\renewcommand{\tH}{\widetilde H}
\providecommand{\tZ}{}\renewcommand{\tZ}{\widetilde Z}
\providecommand{\tU}{}\renewcommand{\tU}{\widetilde U}
\providecommand{\tV}{}\renewcommand{\tV}{\widetilde V}
\providecommand{\tth}{}\renewcommand{\tth}{\widetilde\theta}
\providecommand{\Ocal}{}\renewcommand{\Ocal}{\mathcal O}
\providecommand{\Pfr}{}\renewcommand{\Pfr}{\mathfrak P}
\providecommand{\pfr}{}\renewcommand{\pfr}{\mathfrak p}
\providecommand{\Reg}{}\renewcommand{\Reg}{\operatorname{Reg}}
\providecommand{\Heis}{}\renewcommand{\Heis}{\mathrm{Heis}}
\providecommand{\Cubes}{}\renewcommand{\Cubes}{\mathrm{Cubes}}
\providecommand{\Good}{}\renewcommand{\Good}{\mathrm{Good}}
\providecommand{\Ptag}{}\renewcommand{\Ptag}{\textnormal{\textbf{(P)}}}
\providecommand{\Pctag}{}\renewcommand{\Pctag}{\textnormal{\textbf{(P$^{\mathrm c}$)}}}
\providecommand{\Etag}{}\renewcommand{\Etag}{\textnormal{\textbf{(E)}}}
\providecommand{\Otag}{}\renewcommand{\Otag}{\textnormal{\textbf{(O)}}}
\providecommand{\Ctag}{}\renewcommand{\Ctag}{\textnormal{\textbf{(C)}}}
\section{\texorpdfstring{The wild place in degree two: toral positivity and the nonabelian gluing problem (Chapter~\ref{ch:W})}{The wild place in degree two: toral positivity and the nonabelian gluing problem (Chapter 8)}}\label{W:sec:summary}
Throughout, $F=\Q_3(\zeta_3)$ unless said otherwise, $D$ is the
quaternion division algebra over $F$, $H=D^\times$, $\tH$ is the
$3$-fold scalar-central cover of Chapter~\ref{ch:I}, and $L_D$ is the division-side
character transform of Chapter~\ref{ch:III}, normalized in
Section~\ref{W:sec:transform}.  The labels are those of
Section~\ref{W:sec:intro}.

\begin{enumerate}[leftmargin=*]
\item \textbf{Toral positivity} (Theorem~\ref{W:thm:main}, \Ptag).  For
every quadratic $E/F$ with $p$ odd and every odd $n$ with $\mu_n\subset F$,
the discriminant ratio $\Delta(t)/\Delta(t^n)$ is the character of a
finite permutation representation of $E^\times/F^\times$.  In the running
example the genuine torus multiplicities are $(3,2,2,2)$ for the
unramified torus and $(3,2,1,1,1,1)$ for either ramified torus
(Proposition~\ref{W:prop:all-three-tori}), and integrality of the whole
character matrix is equivalent to positive definiteness of every row
(Proposition~\ref{W:prop:positive}).
\item \textbf{Target A} (Theorem~\ref{W:thm:main-targetA}, \Pctag).
$L_D\mathbf1$ is the character of a single genuine irreducible
representation $\tau$ of $\tH$ of dimension $27$, and
$\tau|_{\tU_1}=\Ind_{\tU_3}^{\tU_1}\lambda_0$ for the explicit character
$\lambda_0$ of \eqref{W:eq:lambda0}.  The proof has three layers: a
Brauer reduction to the pro-$3$ Sylow subgroup
(Theorem~\ref{W:thm:reduction}); a finite design problem on
$V=Q_3/Q_6\simeq\Fq_3^5$ with a unique solution up to sign
(Proposition~\ref{W:prop:design}); and hypothesis (H1), that the
canonical star section is a homomorphism on $K^6$, proved by a
Hilbert-symbol computation in an explicit Kubota--Kazhdan--Patterson
model of the cover (Proposition~\ref{W:prop:H1}).  The label \Pctag\
records that the finite design and the symbol tables are exhaustive
exact computations (Remark~\ref{W:rem:dependencies}); every other step is
a hand proof.
\item \textbf{Target B, reduction and criterion}
(Theorem~\ref{W:thm:reduction-all}, \Ptag; Proposition~\ref{W:prop:criterion}, \Pctag).
For every irreducible $\rho$ with $\theta_\rho|_{\mu_3}=1$,
$L_D\Theta_\rho$ is an irreducible genuine character if and only if its
restriction to $\tU_1$ is a virtual character, and
$L_D\Theta_\rho=\Theta_\tau\cdot F_\rho$ with $F_\rho$ the cube-root
transform of $\Theta_\rho$ (Lemma~\ref{W:lem:product}).
\item \textbf{Target B at levels $\le2$}
(Theorems~\ref{W:thm:levelone}, \ref{W:thm:leveltwo}, \Pctag).  For every
irreducible $\rho$ trivial on $K^3$ (and every norm twist of one),
$L_D\Theta_\rho$ is the character of a genuine irreducible
representation $\tau_\rho$ of dimension $27\dim\rho$, and
$\tau_\rho|_{\tU_1}$ is an explicit sum of irreducible representations
induced from $\tU_3$: nine of dimension $27$ indexed by the characters
of $K^1/K^2$ at level $\le1$, and two of dimension $81$ for the
six-dimensional representations of level two, where the proof rests on
Lemma~\ref{W:lem:Hstructure} and on the freeness of the twisted action of
$Q_1/Q_3$ on the genuine characters of $\hat H$
(Proposition~\ref{W:prop:free}).  Levels $\ge3$ are open
(Sections~\ref{W:subsec:cuberoot}, \ref{W:sec:future}).
\item \textbf{Target C, uniform ingredients}
(Theorems~\ref{W:thm:pthpower}, \ref{W:thm:H1-gen}, \Ptag).  For every odd
prime $p$ and every $F\supset\mu_p$ with $n=p$: the $p$-th power identity
$(w,0)^p=(w^p,0)$ in the Kubota model, and the multiplicativity of the
canonical star section on $F^\times K^{2d}$, $d=pe/(p-1)$.  The Brauer
reduction is expected to be uniform; the finite design is not
(Remark~\ref{W:rem:uniform}).
\item \textbf{Reformulation} (Remark~\ref{W:rem:howe-duality}).  Target~B
for a central character $\theta$ is equivalent to the statement that,
for every level $m$, the explicit kernel
$K^{(m)}_\theta=\Theta_\tau\cdot\tth\tth_1^{-1}\cdot\Delta^{(m)}_\theta$ on
$\tH\times H/K^{m+1}$ is the character of a representation, which then
has the multiplicity-one property of a Howe duality.
\end{enumerate}
What remains, and why we stop at level two, is the subject of
Section~\ref{W:sec:future}.
\endgroup
